\section{设计原则}

\begin{proposition}[最便宜可控变量原则]
当多个障碍项同时升高时，第一项干预应当瞄准那个能够实质改变转变概率、且成本最低的可控变量，而不是瞄准那个在道德上看起来最令人敬佩的变量。
\end{proposition}

\begin{discussion}
这一原则遵循控制设计的逻辑。在反复出现的失败中，提高内在决心或许是可能的，但这种做法通常代价高、噪声大，而且很难验证。相比之下，把手机移远、预先标出下一道题，或者把任务绑定到一个已经反复出现的边界，都会更直接、更低成本地改变状态转变条件。用动力系统语言说，如果局部操作改变了系统在敏感进入点附近的条件，那么小规模操纵也可能非常关键。
\end{discussion}

\begin{proposition}[历史依赖强化原则]
评价一个干预时，不仅要看它是否成功了一次，还要看它是否通过累积的局部历史，让下一次尝试变得更容易。
\end{proposition}

\begin{discussion}
这正是涌现进入实践的地方。如果读者反复使用同一地点、同一启动线索与同一个最小可行进入动作，那么未来的转变就会继承一条被部分训练过的路径。这还不是严格数学意义上已经形成完整吸引子的证明，但它是一个合理的设计推断，来自关于亚稳态、因果涌现与路径依赖的中等证据车道观念：重复的局部结构可以生成更稳定的高层模式。
\end{discussion}